% Enrique Garcia Rivera | Resume for Georgia Tech M.S. Aerospace Engineering
% Final review copy, November 12, 2026. Compile with pdfLaTeX (Overleaf works).
\documentclass[11pt, letterpaper]{article}
\newcommand{\myphone}{(816) 749-9954}
\usepackage[margin=0.5in]{geometry}
\usepackage[utf8]{inputenc}
\usepackage[T1]{fontenc}
\usepackage[default,lining]{raleway} % lining numerals: dates and numbers read cleanly
\usepackage{anyfontsize}
\usepackage{xcolor}
\usepackage{enumitem}
\usepackage{titlesec}

% Accessibility and ATS: real, copyable text in the PDF
\input{glyphtounicode}
\pdfgentounicode=1

\setlength{\parindent}{0pt}
\hyphenpenalty=10000
\exhyphenpenalty=10000
\sloppy
\pagestyle{empty}
\setlist[itemize]{noitemsep, topsep=1pt, leftmargin=1.15em, itemsep=1.5pt, parsep=0pt, before=\raggedright, after=\vspace{3pt}}
\definecolor{darkslate}{RGB}{20,25,35}
\definecolor{accentblue}{RGB}{0,56,168}
\usepackage[colorlinks=true, urlcolor=accentblue, linkcolor=accentblue,
  pdftitle={Enrique Garcia Rivera Resume},
  pdfauthor={Enrique Garcia Rivera},
  pdfsubject={Resume},
  pdflang={en-US}]{hyperref}
\titleformat{\section}{\normalsize\bfseries\color{darkslate}\uppercase}{}{0em}{}[\vspace{0.5pt}\titlerule]
\titlespacing*{\section}{0pt}{8pt}{3pt}

% \role{Title}{Organization}{City, ST}{Dates}
% Dates stay right-aligned with a guaranteed gap; they drop to the next line only if the line is full.
\newcommand{\role}[4]{{\noindent\parfillskip=0pt \textbf{#1}, #2, #3\unskip\nobreak\hfil\penalty50\hskip1.5em\hbox{}\nobreak\hfil #4\par}}

% Portfolio link
\newcommand{\website}{\href{https://enrique-gr.github.io/portfolio/}{enrique-gr.github.io/portfolio}}

\begin{document}

\begin{center}
{\fontsize{18pt}{22pt}\selectfont\bfseries ENRIQUE GARCIA RIVERA} \\[4pt]
\small
\href{tel:+18167499954}{\myphone} \ $|$ \
\href{mailto:e.garciarivera@wustl.edu}{e.garciarivera@wustl.edu} \ $|$ \
\href{https://www.linkedin.com/in/enrique-gr}{linkedin.com/in/enrique-gr} \ $|$ \
\website
\end{center}

\section{Education}
\textbf{Washington University in St.\ Louis} \hfill St.\ Louis, MO \\
Bachelor of Science in Mechanical Engineering \ $|$ \ GPA: 3.65 / 4.00 \hfill Expected May 2027 \\
\textbf{Relevant Coursework:} Flight Dynamics and Aircraft Design (graduate, in progress); Modeling, Simulation, and Control; Numerical Methods and Matrix Algebra; Vibrations; Fluid Mechanics

\section{Research Experience}
\role{UAS Research Intern}{Missouri Institute for Defense and Energy}{Kansas City, MO}{Summer 2026}
\begin{itemize}
\item Built a MATLAB pipeline that batch-processes ArduPilot flight logs from 3 tactical fixed-wing Group 2 UAS, automatically detecting 125 flights in one month and organizing flight activity by aircraft and day.
\item Developed an interactive MATLAB center of gravity model from measured landing-gear weights to evaluate payload layouts and solve for the nose ballast needed to hold the design CG.
\item Rebuilt quadcopters from component-level hardware for a field-experimentation fleet, configuring ArduPilot over MAVLink and DroneCAN with per-vehicle parameter baselines and calibration.
% [CONFIRM before uncommenting] \item Built a DroneCAN servo command chain with a self-leveling control loop [what it leveled; your exact role].
\item Supported UAS flight testing in the field, including ground station setup and composite airframe repair.
\end{itemize}

\role{Research Intern}{Multiplatform Interactive Robotics Laboratory}{Kansas City, MO}{Summer 2025}
\begin{itemize}
\item Built a Teensy 4.0 embedded sensing platform synchronizing two EMG and two IMU streams at 1\,kHz, and ran timed trials across rest, contraction, and motion conditions.
\item Processed EMG and IMU signals in Python (filtering, RMS envelopes, FFT); contributed to an IEEE Body Sensor Networks 2025 abstract, \textit{Toward Electromechanical Measurement of Muscle Activity}.
\end{itemize}

\section{Engineering Experience}
\role{Founding Member and Lead Avionics Engineer}{WashU VTOL}{St.\ Louis, MO}{Sep.\ 2025 -- Present}
\begin{itemize}
\item Lead avionics and flight control integration for the team's Vertical Flight Society Design-Build-Vertical Flight aircraft: a hexacopter prototype (2025--26) and the current 2026--27 quadcopter.
\item Integrated the hexacopter's flight control and sensing system around a Pixhawk 6C running PX4 with isolated 4S/6S power; configured flight modes and failsafes and tuned PID gains.
\item Traced barometer and magnetometer dropouts to motor-induced electromagnetic interference on the I2C bus and restored reliable sensor data by rerouting lines, twisting signal pairs, and adding pull-ups.
\item Validated the integrated avionics in stages: bench, hardware-in-the-loop simulation, and first hover.
% [CONFIRM before uncommenting] \item Defining the avionics architecture for the 2026--27 quadcopter [state only what is decided or built].
\end{itemize}

\role{Structures and Test Engineer}{WashU Design/Build/Fly}{St.\ Louis, MO}{Aug.\ 2024 -- Present}
\begin{itemize}
\item Sized primary wing structure for AIAA Design/Build/Fly and SAE Aero Design aircraft under 2.5\,g loads using ANSYS finite element analysis, iterating on carbon-fiber spar layup.
\item Designed and fabricated mission payload systems, including a fuel-carrying composite pylon and a winch deployment system, and built a drop-test stand to verify payload protection.
\end{itemize}

\section{Teaching}
\noindent\textbf{Undergraduate Teaching Fellow}, Washington University in St.\ Louis \hfill Aug.\ 2025 -- Present \par
\begin{itemize}
\item Hold office hours and grade for Structural Dynamics and Vibrations (Fall 2026) and Electrical and Electronic Circuits (Fall 2025), helping students connect theory to physical systems.
\end{itemize}

\section{Technical Skills}
{\raggedright
\textbf{Dynamics, Control, and Analysis:} MATLAB/Simulink, Python, C++, dynamics modeling, feedback control, PID tuning, signal processing, ANSYS FEA, SolidWorks \\
\textbf{Flight Systems and Test:} PX4, ArduPilot, MAVLink, DroneCAN, Pixhawk, HIL testing, embedded systems, data acquisition, sensor calibration, composites\par}

\end{document}